\documentclass[11pt,a4paper]{article}
\usepackage[T1]{fontenc}
\usepackage{lmodern,amsmath,amssymb,booktabs,tabularx,enumitem}
\usepackage[margin=22mm,headheight=15pt]{geometry}
\usepackage{xcolor,microtype,fancyhdr,hyperref}
\definecolor{navy}{HTML}{17365D}
\hypersetup{colorlinks=true,linkcolor=navy,urlcolor=navy,
pdftitle={Focused Statistical Learning: Decision Trees and Random Forests}}
\pagestyle{fancy}\fancyhf{}
\fancyhead[L]{\small Statistical Learning: essentials}
\fancyhead[R]{\small Trees and random forests}
\fancyfoot[C]{\thepage}
\setlength{\parindent}{0pt}\setlength{\parskip}{.45em}
\setlength{\emergencystretch}{2em}
\setlist{itemsep=.2em,topsep=.3em}
\newcommand{\code}[1]{\texttt{\small #1}}
\newcommand{\idea}[1]{\textcolor{navy}{\textbf{Remember.}} #1}
\begin{document}
{\LARGE\color{navy}Decision trees and random forests}\par
{\large A focused guide to the formulas used in applications A, B, and C}

\section*{1. A tree is a sequence of questions}
A tree routes an input through tests such as ``queue length $q\le3$?'' and
``payload $s\le5$ MB?''. Each terminal node, or \emph{leaf}, stores a prediction
learned from the training rows that reach it. Depth counts tests on the longest
path; a root-only tree has depth zero. A stump has one split and two leaves.

A tree can represent the conditional mean in application A:
\[
f(s,q)=\begin{cases}
40\text{ ms},&q\le3,\\
60\text{ ms},&q>3,\ s\le5,\\
90\text{ ms},&q>3,\ s>5.
\end{cases}
\]
Payload matters only when the queue is long: the tree expresses an interaction
without requiring a product feature. Ordinary numerical tests create axis-aligned
regions and constant predictions within each region. They do not smoothly
extrapolate beyond the observed feature range.

\textbf{Regression: predict the mean inside a leaf.}
For rows $A$ in a leaf, minimizing $\sum_{i\in A}(y_i-c)^2$ gives
\[
\boxed{c_A=\bar y_A,\qquad
\operatorname{RSS}(A)=\sum_{i\in A}(y_i-\bar y_A)^2.}
\]
Indeed the derivative with respect to $c$ is $2|A|c-2\sum_{i\in A}y_i$.
For a feature $j$ and threshold $t$, split into $A_L=\{i:x_{ij}\le t\}$
and $A_R=\{i:x_{ij}>t\}$. Choose the admissible split with smallest
$\operatorname{RSS}(A_L)+\operatorname{RSS}(A_R)$. Reject empty or too-small children.

\textbf{A split you can verify.} Take $x=(1,2,3,4,5,6)$ and
$y=(2,3,2,8,9,10)$. Parent RSS is $208/3$.
\begin{center}
\begin{tabular}{@{}lrrrrr@{}}\toprule
Threshold & 1.5 & 2.5 & 3.5 & 4.5 & 5.5\\
Child RSS total & 53.2 & 39.25 & $8/3$ & 25.25 & 46.8\\\bottomrule
\end{tabular}
\end{center}
The winning split is $x\le3.5$. The leaf predictions are $7/3$ and 9;
training RMSE is $\sqrt{(8/3)/6}=2/3$. The RSS reduction also equals
\[
\boxed{\frac{n_Ln_R}{n_L+n_R}(\bar y_L-\bar y_R)^2,}
\]
so a split helps when the child means differ, with their sizes taken into account.

\textbf{Efficient search.} Sort each feature; test midpoints only between
distinct values. Prefix sums $S_k=\sum_{i\le k}y_{(i)}$ and
$Q_k=\sum_{i\le k}y_{(i)}^2$ give left RSS $Q_k-S_k^2/k$ and the analogous
right RSS. Center responses first to reduce cancellation from a large common
offset. Sorting costs $O(n\log n)$; scanning then costs $O(n)$ per feature.

\idea{Split selection fits the training data. A low training RSS is not an
estimate of how well the tree predicts new observations.}

\newpage
\section*{2. Classification: keep the leaf probabilities}
For a binary target $y\in\{0,1\}$, a leaf estimates the positive probability as
\[
\boxed{\widehat\pi_A=\frac{\#\{i\in A:y_i=1\}}{|A|}.}
\]
Classify as 1 if $\widehat\pi_A\ge t$ under our explicit tie convention.
The threshold $t$ is a decision choice; it is not part of learning the split
criterion. A leaf with one positive row gives probability 1 but very little
evidence about future certainty. Minimum leaf size helps control this variability.

\textbf{Gini impurity measures class mixing.}
Write $G(A)$ for the Gini impurity of node $A$. For class fractions $p_1,\ldots,p_K$,
\[
G(A)=1-\sum_{k=1}^Kp_k^2,\qquad
G(A)=2\widehat\pi_A(1-\widehat\pi_A)\quad\text{in the binary case}.
\]
Gini is zero for a pure leaf and largest when classes are balanced. It is the
probability that two independent draws from the leaf's empirical label distribution
disagree. Entropy $H(A)=-\sum_kp_k\log p_k$, with $0\log0=0$, is another
common measure; it also favors purer groups.

\textbf{Weight child impurities by their sizes.}
Here $I(A)$ is a \emph{generic name for the chosen impurity}, not a new measure:
use $I(A)=G(A)$ for Gini, or $I(A)=H(A)$ for entropy.
With $n_A=|A|$ and child sizes $n_L,n_R$, maximize
\[
\boxed{\operatorname{gain}=I(A)-\frac{n_L}{n_A}I(A_L)
-\frac{n_R}{n_A}I(A_R).}
\]
Without these weights, a tiny pure child would receive too much influence.
For regression, $I(A)=\operatorname{RSS}(A)/n_A$ gives the same split ranking
as minimizing total child RSS.

\textbf{A better probability split can leave accuracy unchanged.}
\begin{center}
\begin{tabular}{@{}lrrrr@{}}\toprule
Node & Positive & Negative & $\widehat\pi$ & Gini\\\midrule
Parent & 8 & 2 & 0.8 & 0.32\\
Left child & 5 & 0 & 1 & 0\\
Right child & 3 & 2 & 0.6 & 0.48\\\bottomrule
\end{tabular}
\end{center}
Weighted child Gini is $0.5(0)+0.5(0.48)=0.24$, so gain is $0.08$.
At threshold 0.5 both children still predict 1, leaving the same two mistakes.
Gini detects a useful change in probability structure that immediate accuracy misses.

\textbf{Categorical inputs need explicit treatment.}
In Carseats, shelf location is nominal. One-hot indicators such as
\code{ShelveLoc\_Good} permit the test ``Good versus not Good''.
Arbitrary integer codes silently impose a split ordering. Fit the encoder on
the fitting rows and specify what happens to a category not seen there.
The notebooks use all-zero indicators for an unknown category.

The target is $Y=\mathbf1\{\code{Sales}>8\}$, so \code{Sales} must be
excluded from the inputs. A predictor measured only after the outcome would also
be unusable for a genuine advance prediction.

\idea{Use leaf probabilities for probability metrics and cost-sensitive decisions.
A hard class prediction throws away their degree of certainty.}

\newpage
\section*{3. Control a tree's complexity}
Tree construction is greedy: choose the best current split and repeat in each
child. The first split is not reconsidered after later splits. A useful deeper
structure may require an initially weak split, so greedy fitting need not find
the best possible tree. Small changes in training rows can also change early tests.

\textbf{Two direct controls.} Smaller maximum depth limits the number of tests;
larger minimum leaf size requires each prediction to average more observations.
More complex trees usually fit training data better but may fit noise.
Strictly increasing feature transformations preserve candidate partitions in
exact arithmetic, so trees usually need no standardization. One-hot encoding
and feature availability still matter.

\textbf{Pruning puts a price on every leaf.}
Grow a larger tree and consider its subtrees. Minimize
\[
\boxed{R_\alpha(T)=R(T)+\alpha L(T),\qquad \alpha\ge0,}
\]
where $L(T)$ is the leaf count. In regression,
$R(T)=\sum_{\ell}\operatorname{RSS}(A_\ell)/n$; for Gini classification,
$R(T)=\sum_\ell(n_\ell/n)G(A_\ell)$. Larger $\alpha$ favors fewer leaves.
The regression penalty has squared-response units; the Gini penalty has a
different scale. Using raw RSS instead of normalized RSS multiplies equivalent
penalties by $n$.

For the six-row example, consider a tree with a root test at 3.5 and an
additional right-branch test at 4.5:
\begin{center}
\begin{tabular}{@{}lrrl@{}}\toprule
Subtree & Leaves & $R(T)$ & Penalized risk\\\midrule
Keep both tests & 3 & $7/36$ & $7/36+3\alpha$\\
Keep root only & 2 & $4/9$ & $4/9+2\alpha$\\
No tests & 1 & $104/9$ & $104/9+\alpha$\\\bottomrule
\end{tabular}
\end{center}
The first two options tie at $\alpha=1/4$; the last two at $100/9$.
Choose the smaller tree at a tie. These are training tradeoffs among subtrees;
cross-validation is still needed to select predictive complexity.

\textbf{The same evaluation discipline as regression applies.}
\begin{enumerate}
\item Make 60/20/20 train/validation/test partitions before target exploration;
      stratify binary labels in these independent-data examples.
\item Compare a prespecified parameter grid using five-fold training CV.
      Grow a fresh tree and fit preprocessing inside each fold.
\item Select by RMSE for regression or Brier score for binary probabilities.
      Fit the selected configuration on the training partition.
\item Use validation for an audit, probability/threshold checks, or exploratory
      importance. Freeze the model and policy before the final test.
\end{enumerate}
Do not grow a full-data tree first and reuse its learned splits across CV folds.
Repeated customers or future-time prediction need appropriate group/time splits.

\idea{Pruning changes the tree. Cross-validation chooses the setting. A separate
test estimates the performance of the resulting frozen procedure.}

\newpage
\section*{4. Bagging and random forests: average unstable trees}
\textbf{Bagging.} Draw $n$ training indices \emph{with replacement}, fit a tree,
and repeat $B$ times. Repeated indices retain their multiplicity. Average the
predictions:
\[
\widehat f_B(x)=\frac1B\sum_{b=1}^B\widehat f_b(x),\qquad
\widehat\pi_B(x)=\frac1B\sum_{b=1}^B\widehat\pi_b(x).
\]
Regression averages numerical responses; classification here averages positive
probabilities. Always check class-column ordering when reading probability arrays.

\textbf{Averaging probabilities is not the same as majority voting.}
For five tree probabilities $(0.2,0.3,0.4,0.9,0.95)$, the average is 0.55,
so the averaged-probability rule predicts 1 at threshold 0.5. Thresholding each
tree first gives three votes for 0 and two for 1: majority voting predicts 0.
Record which rule the implementation uses.

\textbf{Why averaging helps.} At one fixed input, suppose each fitted tree
prediction has variance $\sigma_T^2$ across repeated training experiments,
and pairs have common correlation $\rho$. Then
\[
\boxed{\operatorname{Var}\!\left(\frac1B\sum_b\widehat f_b\right)
=\sigma_T^2\left(\rho+\frac{1-\rho}{B}\right).}
\]
Expand the variance of a sum: $B$ individual variance terms plus $B(B-1)$
covariance terms, all divided by $B^2$. With variance 16 and correlation 0.25,
ensemble variance is 16, 5.2, and 4.12 for $B=1,10,100$; the limit is 4.
Adding trees reduces the unshared variation but cannot remove the shared part.

\textbf{Random forests also randomize candidate features.}
At every node, offer a random subset of $m$ columns to the split search rather
than always offering all $p$. This can make trees less similar. It can also
weaken individual splits, so choose $m$ by predictive evaluation.
With $m=p$, bootstrap forests behave as bagged trees under the same controls.
The subset is sampled anew at each node, not once for the whole forest.

\textbf{Two controls with different jobs.}
Increasing $B$ stabilizes the average and increases time/storage. Changing
depth, minimum leaf size, or $m$ changes the fitted rule and its bias/variance
tradeoff. Many trees cannot rescue unavailable features, leakage, a wrong target,
or irreducible outcome noise.

\textbf{Do not confuse two notions of correlation.}
Correlating two trees' prediction vectors across many inputs measures how
similarly their rules vary across those inputs. It is not the fixed-input,
repeated-training correlation $\rho$ in the formula above. Application C reports
that descriptive similarity and separately resamples training data to study
prediction variability at fixed inputs.

\idea{Randomization makes diversity possible; averaging makes diversity useful.
Whether it helps this prediction problem is an empirical question.}

\newpage
\section*{5. OOB evaluation, importance, and decisions}
\textbf{A bootstrap leaves some training rows out.}
The probability that one row is absent from $n$ draws is
\[
q_n=(1-1/n)^n\approx e^{-1}\approx0.368.
\]
Its expected OOB tree count is $Bq_n$. For row $i$, average only over the
trees that did not use it:
\[
\boxed{\widehat\pi_i^{\rm OOB}
=\frac1{|\mathcal B_i|}\sum_{b\in\mathcal B_i}\widehat\pi_b(x_i),\qquad
\mathcal B_i=\{b:i\text{ absent from bootstrap }b\}.}
\]
If $|\mathcal B_i|=0$, its OOB prediction is undefined. Exclude that row from
the current OOB metric and report coverage; never replace it by zero.
With few trees, the coverage and probability estimates can be unstable.
OOB is an internal assessment, not an independent final test after extensive
tuning. Learned full-data preprocessing can also compromise its interpretation.

\textbf{Feature importance answers a specific question.}
Mean decrease in impurity (MDI) totals training split gains credited to a feature,
weighted by the fraction of training mass at each node, then normalizes across
features. It can favor features with many candidate thresholds and reflect overfit.

For validation permutation importance under Brier loss, use
\[
\operatorname{PI}_j=
\operatorname{Brier}(y_{\rm val},\widehat\pi(X_{\rm val}^{\rm perm(j)}))
-\operatorname{Brier}(y_{\rm val},\widehat\pi(X_{\rm val})).
\]
Permute one column, predict without refitting, and repeat the permutation to
measure Monte Carlo variability. Positive importance means the shuffle hurt.
Correlated predictors can substitute for each other; shuffling one also creates
unusual combinations. Jointly shuffle a redundant pair with one shared row
permutation to assess reliance on the pair while preserving their relationship.
Neither method proves causation; MDI and permutation loss have different scales.

\textbf{Evaluate probability quality and decision quality separately.}
\[
\operatorname{Brier}=\frac1m\sum_i(y_i-\widehat\pi_i)^2,\qquad
\text{cost per row}=\frac{c_{\rm FP}\mathrm{FP}+c_{\rm FN}\mathrm{FN}}m.
\]
Pure leaves can give exactly 0 or 1. A wrong probability-zero event has infinite
theoretical log loss. The notebooks explicitly clip to $[10^{-12},1-10^{-12}]$
for reported log loss, applying the same rule to every candidate.

With appropriate probabilities and zero cost for correct decisions, the
cost-minimizing threshold is $c_{\rm FP}/(c_{\rm FP}+c_{\rm FN})$.
Carseats uses costs 1 and 3, giving $t=0.25$. Its notebook also selects a
validation threshold, then freezes it. Report confusion counts, precision,
recall, and cost rather than accuracy alone.

\idea{Training importance, internal OOB scores, and final test performance are
different quantities. State which data and which loss produced each result.}

\newpage
\section*{6. The three companion notebooks}
Each notebook runs independently and places task statements before readable code.
Simulation seed is 17 and split seed is 23. Tables and figures are saved after
execution. Full task details and conventions are in the notebooks.

\subsection*{A. A service with regime changes}
\code{02\_A\_service\_regimes.ipynb}
\[
s\sim U(0,10),\quad q\sim\operatorname{Poisson}(3),\quad
\epsilon\sim N(0,8^2),\qquad
y=40+20\mathbf1\{q>3\}+30\mathbf1\{q>3,s>5\}+\epsilon.
\]
Generate 1,200 independent requests. Verify a prefix-sum stump against every
candidate in the six-row example. Compare a mean baseline, a linear model,
a stump, and a depth-two tree through their prediction surfaces. Select pruning
strength by training CV, then bootstrap training rows to inspect split and
prediction stability. Finish with frozen-model test errors and complexity/runtime.

\textbf{Main lesson:} the tree family matches this piecewise-constant simulator.
That explains its advantage here; it does not establish a universal ranking.

\subsection*{B. Carseats: categorical features and probabilities}
\code{02\_B\_carseats.ipynb}

Use the local 400-row \code{Carseats.csv}; define high sales as
$Y=\mathbf1\{\code{Sales}>8\}$, giving 164 positives. Exclude sales from
inputs. Fit one-hot encoding within folds, compare a logistic baseline, select
depth/leaf controls and a pruning penalty using training CV Brier score, and
draw a shallow tree with leaf counts and positive fractions.

Use validation to audit probabilities and select an alarm threshold for costs
1 and 3. Compare the frozen policy with prespecified probability baselines
on the final test. Translate two paths back into ordinary feature statements.

\textbf{Main lesson:} a tree predicts probabilities from local counts; small
leaves and a small evaluation sample demand careful interpretation.

\subsection*{C. Bagging, decorrelation, OOB, and importance}
\code{02\_C\_forest\_diagnostics.ipynb}

Generate 2,000 rows with 20 columns: $x_1,x_3,x_5,\ldots,x_{20}$ are
independent $U(-1,1)$ draws, $x_2=x_1+N(0,0.03^2)$, and
$x_4\sim\operatorname{Bernoulli}(0.5)$ is independent. Labels follow
\[
Y\mid x\sim\operatorname{Bernoulli}(\pi(x)),\qquad
\pi(x)=\sigma\bigl(-0.7+2.5x_1+2\mathbf1\{x_3>0.2\}\bigr).
\]
Compare bagging and forests with 300 trees; select $m$ and minimum leaf size
by training CV. Independently build bootstrap trees and audit row-specific
OOB averages, coverage, and convergence through 600 trees. Compare tree-rule
similarity, resampled prediction variability, MDI, and validation permutation
importance, including a joint permutation of $(x_1,x_2)$.

\textbf{Main lesson:} averaging, internal evaluation, and interpretation each
require their own checks. Keep generating probabilities out of the input columns.

\small\textbf{Further reading.} The full second lecture is in \code{PdfLectures}.
The bundled book is James et al., \emph{An Introduction to Statistical Learning
with Applications in R}, 2nd ed. (2021), Chapter 8.1, 8.2.1, and 8.2.2;
Chapter 5 supplies bootstrap and validation background.
\end{document}
